\documentclass[11pt]{article}
\usepackage{amsmath,amssymb,amsthm,fullpage}
\newtheorem{proposition}{Proposition}
\title{Orthants at Coincident Reciprocal Branches}
\author{}
\date{November 23, 2026}
\begin{document}
\maketitle

\paragraph{Status.}
This note studies first-order perturbations of the exact full-cube
identities in \texttt{reciprocal-histogram-fibers.tex}. It proves a
finite seven-point interpolation statement and a scoped capacity bound
for these tangent families. Neither accepted asymptotic target is
proved. No model definition is changed. The arguments are written
proofs, with exact rational calculations supplied by
\texttt{reciprocal-histogram-tangents-check.py}; they are not Lean proofs.

\section{A finite orthant gain from two full-cube branches}
On the five-dimensional Boolean cube set
\[
 a=x_1-x_2,\qquad b=x_2-x_3,\qquad
 D=1+4(x_1+x_2+x_3)+8x_4+16x_5.
\]
For $h=0,1,2$ and a permutation $\pi$ of $\{0,1,2\}$ put
$B_{\pi,h}=D+ha+\pi(h)b$.
Take the two permutations
\[
 \pi=(0,2,1),\qquad \tau=(1,0,2).
\]
For every Boolean input the two denominator triples have the same
multiset. Indeed one of $a=0$, $b=0$, or $a+b=0$ holds. In the first
two cases this follows from the permutation property; in the third,
both difference multisets $\{h-\pi(h)\}_h$ and
$\{h-\tau(h)\}_h$ equal $\{-1,0,1\}$.
Every denominator has constant coefficient one and coordinate slopes
\[
 4+h,\quad4+\pi(h)-h,\quad4-\pi(h),\quad8,\quad16,
\]
all strictly positive. Consequently
\[
 S_\pi(x)=\sum_{h=0}^2B_{\pi,h}(x)^{-1}
          =\sum_{h=0}^2B_{\tau,h}(x)^{-1}=S_\tau(x)
\]
is an exact identity between admissible original-model sums.

Perturb a branch by replacing its denominators with
$B_{\pi,h}+\epsilon(d_h+p_hb)$. The derivative at zero is
\[
 T_{\pi,d,p}(x)=-\sum_{h=0}^2
            \frac{d_h+p_hb(x)}{B_{\pi,h}(x)^2}.
 \tag{1}\label{eq:tangent}
\]
The direction $(d,p)$ has six independent real coordinates. Let $E$
be the following ordered seven-point subset of the original cube:
\[
\begin{split}
 &(1,0,1,1,0),\ (1,0,1,1,1),\ (1,0,0,1,0),\ (0,1,1,0,0),\\
 &(1,1,0,0,0),\ (1,1,0,1,1),\ (0,1,0,1,1).
\end{split}
\]
\begin{proposition}
Every strict sign table on $E$ is realized by a sum of six affine
fractions, all with strictly positive coordinate slopes and positive
denominators on the entire five-cube.
\end{proposition}
\begin{proof}
For each of the two branches form the $7$ by $6$ matrix whose row at
$x\in E$ is
\[
 \bigl(-B_{\pi,0}^{-2},-B_{\pi,1}^{-2},-B_{\pi,2}^{-2},
       -bB_{\pi,0}^{-2},-bB_{\pi,1}^{-2},-bB_{\pi,2}^{-2}\bigr).
\]
Its entries are explicitly specified rational numbers. The companion
checker calculates all seven signed maximal minors, verifies they are
nonzero, and verifies that their sign vectors for the two branches
are, respectively,
\[
 (-,+,-,-,-,+,+),\qquad(-,+,-,-,-,+,-).
\]
They are neither equal nor opposite. It also verifies the corresponding
cofactor vectors are orthogonal to every matrix column.

Here is why those finite calculations prove all $128$ strict labels
are obtained. A hyperplane $\nu^\perp\subset\mathbb R^7$, with every
$\nu_i\ne0$, meets every strict orthant except the two with signs
$\operatorname{sgn}\nu$ and $-\operatorname{sgn}\nu$.
For any other sign vector $s$, the numbers $s_i\nu_i$ include both
signs, so positive magnitudes can be chosen to make their weighted
sum zero. This produces a vector in the requested orthant of the
hyperplane. The two branch hyperplanes have disjoint exceptional
pairs, hence their union meets every strict orthant.

For any such tangent witness, subtract the common baseline $S_\pi$
from the perturbed chosen branch. This is a six-fraction score.
Its quotient by $\epsilon>0$ converges at every point of $E$ to
\eqref{eq:tangent}. Strict inequalities and strictly positive
denominator coefficients therefore persist for sufficiently small
$\epsilon$. Thus it realizes the requested signs with original
admissible atoms.

The checker independently enumerates all $128$ tangent labels and
checks actual finite rational perturbations for each of them, including
every denominator at all $32$ Boolean inputs. More explicitly, let
$Q=\max_h(|d_h|+|p_h|)$ and
$\mu=\min_{x\in E}|T_{\pi,d,p}(x)|$. It uses
\[
 \epsilon=\min\left\{\frac1{4(1+Q)},
                  \frac{\mu}{24(1+Q)^2}\right\}>0.
\]
The identity
\[
 \frac1{B+\epsilon u}-\frac1B
 =-\frac{\epsilon u}{B^2}
  +\frac{\epsilon^2u^2}{B^2(B+\epsilon u)}
\]
shows that, after division by $\epsilon$, the total error is at most
$6\epsilon Q^2<\mu/2$, since $B\geq1$ and
$B+\epsilon u>1/2$. The coefficient checks follow from
$\epsilon Q<1/4$ and the original slopes being at least two.
\end{proof}

\paragraph{Exact diagnostic beyond seven points.}
The same script verifies that the two six-dimensional tangent spaces
have a ten-dimensional joint linear span on all $32$ inputs. On the
eight-point set recorded in its JSON output, each individual branch
has exactly $240$ strict orthants, and their union has exactly $244$
of $256$. The enumeration includes exact witnesses and verifies that
the central arrangements are simple, which makes its count complete.
This is one finite configuration. It neither proves nor disproves a
growing interpolation gain.

\section{Full coefficient perturbations on one core face}
The preceding regularity calculation uses only constant and core
coefficient perturbations. Arbitrary numerator and filler-coordinate
perturbations enlarge it considerably. There is nevertheless a useful
bound that applies to all these first-order directions simultaneously.

Fix a core assignment, leave $d$ Boolean filler coordinates
$y=(y_1,\ldots,y_d)$ free, and put
$t(y)=t_0+\sum_iw_iy_i$. Suppose every base denominator on this
face is of the form $t(y)+q_h$, positive at every filler input.
Both the $q_h$ and the selected base branch may vary. A first-order
perturbation of numerator and denominator coefficients at a base
whose numerators are constant has the form
\[
 G(y)=\sum_{h=1}^H
 \left(\frac{A_h(y)}{t(y)+q_h}
              +\frac{C_h(y)}{(t(y)+q_h)^2}\right),
 \tag{2}\label{eq:fulltangent}
\]
where $A_h,C_h$ are arbitrary affine functions of $y$.

\begin{proposition}
If all strict labels on this $d$-cube are realized by scores
\eqref{eq:fulltangent}, allowing arbitrary branches and pole lists,
then
\[
 2Hd+1\geq2^d.
\]
The same conclusion applies to compilers obtained by subtracting the
common base score and taking a nonzero first-order limit.
\end{proposition}
\begin{proof}
Multiply a score by the positive quantity
$P(t)^2=\prod_h(t+q_h)^2$. Its signs are unchanged and its values
belong to the fixed linear space
\[
 V=\operatorname{span}\{t(y)^j,\ y_it(y)^j:
          1\leq i\leq d,\ 0\leq j<2H\}.
\]
The largest degree is $2H-1$, from the simple-pole terms. This
space is independent of all poles and branches. For $1\leq j<2H$,
\[
 t^j=t_0t^{j-1}+\sum_iw_i y_it^{j-1}.
\]
Thus $V$ is spanned by its constant function and the $2Hd$ functions
$y_it^j$, so $\dim V\leq2Hd+1$.

Any proper linear subspace of $\mathbb R^{2^d}$ misses a strict
orthant. Indeed choose a nonzero orthogonal vector $\nu$ and prescribe
the sign of $\nu_x$ at its nonzero coordinates, with arbitrary signs
at the remaining coordinates. Every vector in that orthant has
strictly positive inner product with $\nu$. If $2Hd+1<2^d$, this
argument applied to $V$ supplies a label absent from every branch.
\end{proof}

\paragraph{Scope of this bound.}
It uses the common affine filler statistic at the base and the
first-order form \eqref{eq:fulltangent}. Finite perturbations with
independently varying filler slopes, higher-order cancellations,
and unrestricted original representations are not covered. With
$k$ core triples there are $d=n-3k$ filler bits. The lower bound
is then on the scale $2^n/(n8^k)$, so a growing number of core
triples remains compatible with the requested logarithmic savings.
Consequently this result does not establish the accepted global
lower bound, and the finite seven-point construction does not
establish the accepted global upper bound.
\end{document}
